\section{Variants of the prediction-stability condition}
\label{sec:theory-variants}

\subsection{The tolerance is determined by the margin gap}

The constants in the main condition arise from the gap between the training
and test margins. They need not be treated as separate modeling choices.

\begin{corollary}[General training and test margins]
\label{cor:general-margin-gap}
Fix $m>0$ and $0\le u<m$, and put $\Delta=m-u$.
Replace the training likelihood by $f(X_i)z_w(X_i)>m$, the test error by
$f(X)z_w(X)\le u$, and the clipping map by
$T_m(z)=\max\{-m,\min\{m,z\}\}$.
Require clipped-score variation at most $\Delta/4$ on the calibrated ball
of radius $c/\sqrt{1+HE_v}$. Then the main theorem holds with a sample
bound polynomial also in $m$ and $\Delta^{-1}$.
\end{corollary}
\begin{proof}
Use readout coefficient $a=\max\{2,2m,256c\}$ in the target cylinder.
Its noisy signed output is greater than $m$, its calibrated radius is still
at most $\tau_0$, and its Gaussian negative log-mass remains polynomial
in $m,c,d,L,C$. After the usual readout truncation use
\[
 r=\min\{1,2\Delta/3,c/\sqrt{1+HB}\}.
\]
Successful compression changes the clipped score by at most
$\Delta/4+r/8\le\Delta/3$.
Training scores are therefore at least $m-\Delta/3$, whereas a test
margin failure has candidate score at most $u+\Delta/3$.
Use a $[0,1]$-valued ramp which is zero above the former threshold and
one below the latter; these thresholds are separated by $\Delta/3$.
The same relative-Chernoff argument proves the result.
\end{proof}

The main theorem is $m=1/2$, $u=1/4$, so $\Delta/4=1/16$.
If only the sign is scored at test time, take $u=0$: the permissible
score-variation tolerance increases to $1/8$. This concerns the raw-output
test margin, separately from whether posterior confidence is judged at
$1/2$ or $2/3$.

\subsection{Disturbances confined to valid activation states}

One may restrict the support condition to disturbances whose complete
noisy hidden trajectory lies in $[-1,1]^N$ at every layer. This permits
perturbations only to values within the range of the tanh activation.
The same main theorem holds. Indeed, before grid rounding, every successful
virtual hidden state in the constructive entropy compressor is a tanh
output. Thus the disturbances used in its prediction comparison satisfy
this extra state restriction. The separate output injection and subsequent
grid-rounding error are treated exactly as before. The original target
cylinder was stable even without the state restriction and therefore
retains its probability lower bound.

This condition is weaker than unrestricted additive-disturbance stability.
A negative empirical attack for this variant must verify the state
constraint at every intermediate layer; an attack outside the activation
range does not refute it. The primary empirical condition can remain the
unrestricted one, with this variant stated separately.

\subsection{A condition involving hidden disturbances only}

The scalar output injection can also be omitted from the prior condition.
For $E_v>0$ define the hidden-only calibrated radius
\[
 r^{\rm hid}_c(w)=\frac{c}{\sqrt{H E_v(w)}}.
\]
When $v=0$, the output is identically zero and hidden disturbances cannot
change it, so declare hidden-only score stability automatically satisfied.
Require score variation at most $1/16$ for every hidden disturbance with
$\sum_s\|\eta^s\|^2/N\le(r^{\rm hid}_c(w))^2$, with no output injection.
The main theorem continues to hold for this Gaussian-conditioned prior.

To see this, use the same target cylinder, which has
$r^{\rm hid}_c\le\tau_0$, and after readout truncation choose
\[
 r=\min\{1/30,c/\sqrt{HB}\}.
\]
The raw compression lemma supplies hidden disturbances of norm less than
$r/2$, plus an output injection of magnitude less than $r/2$.
Use stability only for the hidden disturbances and treat the output
injection as an approximation error. Including grid rounding, the total
clipped-score error is at most
\[
 \frac1{16}+\frac r2+\frac r8\le\frac1{12}.
\]
The training/test ramp proof is unchanged. In a nonexpansive stack the
hidden-only gain bound is $\sqrt{H E_v}$, so this calibration yields raw
output variation at most $c$ without the additional scalar-noise term.

\subsection{Amplification is not evidence about unmodified trained weights}

Every finite-cube network with no zero clean outputs eventually satisfies
the calibrated clipped condition after a sufficiently large positive
readout amplification. To verify this, continuity and finiteness of the
input cube give a hidden-disturbance radius $\rho>0$ and $b>0$ such that
the unamplified noisy output has its clean sign and magnitude at least $b$
throughout that radius. Scaling the readout by $A$ makes the calibrated
radius tend to zero while the noisy signed output is at least
$Ab-r_c(w^{(A)})$. Eventually this exceeds $1/2$, so the clean and noisy
clipped scores agree exactly. The readout energy simultaneously grows by
$A^2$. The scale needed may be enormous; this observation gives no
polynomial target-mass bound for arbitrary neural representations and
cannot justify reporting rescaled networks as evidence about ordinary
training.
